\chapter{Fuentes analizadas}

\section{Documentación administrativa y de formato}

Para la redacción de esta memoria se revisaron los documentos incluidos en \texttt{INFO/UTILES}. Esta documentación sirvió para fijar el formato de portada, comprobar los datos administrativos del trabajo y contrastar la descripción inicial del proyecto con el desarrollo realizado. Los documentos más relevantes fueron:

\begin{itemize}[leftmargin=*]
  \item portada y normas básicas de estilo de la Escuela de Ingenierías Industrial, Informática y Aeroespacial.
  \item solicitud de preinscripción del TFG, con título, descripción y tutoría.
  \item resolución favorable de preinscripción de fecha 18 de mayo de 2026.
  \item propuesta \textit{CognitiveCarSimulatorReloaded - IA}.
  \item memoria previa \textit{CognitiveCarSimulator}.
  \item informe de seguimiento de prácticas HP SCDS.
  \item memoria final de prácticas del alumno.
  \item estructura orientativa propuesta por el tutor.
  \item ficha y licencia del asset de semáforo descargado desde Fab.
\end{itemize}

También se incorporaron las especificaciones reales del equipo de desarrollo y pruebas, confirmadas mediante comunicación directa con el tutor de HP: OMEN 45L GT22-0021ns, Intel Core i9-12900K, 32 GB de memoria RAM, unidad SSD de 2 TB y tarjeta gráfica NVIDIA GeForce RTX 3080 Ti. Esta confirmación afecta únicamente a las características técnicas del equipo; el precio de compra, la antigüedad y el consumo eléctrico siguen tratados como estimaciones en el presupuesto.

\section{Actas de reunión}

Se han analizado ocho actas de seguimiento, comprendidas entre el 27/11/2025 y el 19/03/2026. Estas actas permiten justificar las decisiones de arquitectura y los cambios de enfoque realizados durante el proyecto.

\Needspace{8\baselineskip}
\section{Seguimiento en GitLab}

Se revisó el fichero \path{INFO/UTILES/gitlab_issues_final_limpio.csv}, exportado a partir de las issues y milestones del proyecto en GitLab. El CSV contiene 40 issues cerradas agrupadas en 15 sprints, entre el 7 de noviembre de 2025 y el 1 de julio de 2026. Para la memoria se depuraron títulos y descripciones, manteniendo fechas, identificadores, comentarios y enlaces. Esta fuente refuerza la planificación real, la línea evolutiva y la relación entre tareas cerradas y cambios técnicos.

\section{Descripciones del programa}

Las descripciones fechadas de los Blueprints se utilizaron para reconstruir el estado de la implementación sin depender únicamente de la versión abierta en el editor. Para el gestor se revisó la carpeta \texttt{EVOLUTIONMANAGER}, donde se documentan la creación de generaciones, la asignación de modelos, la mutación, la persistencia y el cambio al modo de test. La carpeta \texttt{REDNEURONAL} permitió comprobar la composición de las entradas, las dimensiones de \texttt{WeightsIH} y \texttt{WeightsHO}, la inferencia y las guardas frente a modelos incompatibles.

El cálculo de aptitud y su instrumentación proceden de \texttt{FITNESS}. Se revisaron versiones hasta el 6 de julio de 2026, incluidas las nuevas métricas de genealogía, rotondas, navegación fuera de calzada, solapes entre vehículos, bloqueo en cruces, bonus de espera en cola, cedas liberados y correcciones de escala de distancia frontal y aceleración. Para la señalización se consultó \texttt{INTERSECTIONS}, con la validación aprendida de semáforos y stops, opciones de trayectoria separadas por modo, parejas de dirección y trigger, lógica específica de rotondas, colas dinámicas de liberación en cedas y separación de subtipos de infracción. La actualización del 6 de julio se utilizó únicamente para describir la corrección de stop basada en \texttt{LearnedMinTimeCache}, posición legal y timeout; no se trató como resultado experimental nuevo. Por último, \texttt{OTROS} reúne macros auxiliares, trazas, detección de estancamiento y tareas pendientes que ayudan a interpretar dependencias entre módulos.

La carpeta \path{INFO/FOTOS_PROGRAMA} se revisó también tras su reorganización en subcarpetas de señales, mapa y Blueprints. Las imágenes de señalización ya incorporadas al capítulo técnico proceden de esa fuente. Entre las capturas nuevas de Blueprints se seleccionó únicamente el bloque de lectura sensorial, porque ayuda a entender la arquitectura de percepción sin duplicar la explicación textual de la red ni introducir grafos demasiado extensos.

\Needspace{15\baselineskip}
\section{Resultados experimentales}

Los resultados documentados proceden de \path{INFO/ANALISIS_DATOS/Analisis}. La referencia principal de entrenamiento es \texttt{train\_2026-06-18\_11-21}, con 548 generaciones y 16432 registros de vehículos. La evaluación agrega las cuatro carpetas de test del 18 de junio, que suman 21 ejecuciones y 1134 registros sobre 54 spawns. La marca horaria forma parte del identificador y no se utiliza como variable experimental.

Para reconstruir la evolución se revisaron además los bloques de los días 2, 5, 8, 11, 15, 16, 17, 22, 23, 26, 29 y 30 de junio. Las carpetas \texttt{test\_2026-06-11\_11-55} y \texttt{test\_2026-06-11\_13-08} contienen exactamente el mismo CSV de detalle y se contabilizan una sola vez. El test del día 15 conserva 431 de las 432 filas declaradas. En las sesiones nuevas también se conserva visible la diferencia entre evaluaciones declaradas y registros de detalle cuando existe un pequeño descuadre; en \texttt{test\_2026-06-30\_16-56}, por ejemplo, el analizador selecciona la generación presente en \texttt{Test\_Summary.csv} y excluye filas de una generación sin resumen cerrado. Esta incidencia y las coberturas de cada sesión se mantienen visibles en el capítulo de resultados.

\Needspace{8\baselineskip}
También se analizaron los entrenamientos de los días 11, 16, 17, 18, 22, 23, 25, 26, 29 y 30 de junio, además de los entrenamientos del 1 y 2 de julio. Los bloques posteriores al día 18 no sustituyen a la referencia estable porque muestran regresiones normativas, entrenamientos incompletos o resultados demasiado variables entre sesiones. Se conservan para documentar el efecto de los cambios de fitness, stop, ceda, rotondas, diagnóstico de solapes y bloqueo en cruce. Las sesiones del 3 de julio se trataron aparte como \texttt{FreeRun}, ya que exportan eventos terminales y coches vivos censurados en lugar de generaciones cerradas; la carpeta \texttt{test\_2026-07-03\_15-31} se utilizó como diagnóstico posterior de \texttt{NAV\_VIOLATION}, \texttt{STOP\_VIOLATION\_EXIT}, rotondas y hotspots por trigger. La figura del capítulo 7 se ha regenerado en PDF vectorial a partir del \texttt{Training\_Summary.csv} del día 18, y tanto los datos como el código PGFPlots quedan archivados en \texttt{MEMORIA/assets}.

\clearpage
\section{Conversaciones con asistentes de IA}

La carpeta \path{INFO/CONVERSACIONES_IA} contiene conversaciones mantenidas durante el desarrollo con ChatGPT, Claude y Gemini. Se han revisado 48 ficheros Markdown: 13 de ChatGPT, 28 de Claude y 7 de Gemini. En conjunto reúnen 769 bloques de consulta o intervención del usuario. No se han tratado como fuente primaria de resultados, sino como trazabilidad del proceso de trabajo: preguntas de depuración, análisis de entrenamientos, selección de modelos, organización del proyecto, uso de Unreal Engine y preparación de scripts auxiliares.

El criterio de incorporación ha sido selectivo. Las consultas menores sobre manejo del editor, avisos de LaTeX o dudas puntuales solo se han usado para entender el contexto. En cambio, sí se han tenido en cuenta las conversaciones que ayudan a explicar decisiones técnicas: precisión de parada, detección de \texttt{INVALID\_BRAIN}, coches inmóviles sin justificación, separación train/test, duplicados de modelos, diagnóstico de rotondas, análisis de caídas de rendimiento y regresiones por \texttt{STOP\_VIOLATION}. El anexo B recoge una versión corregida y mejorada de los prompts más representativos, junto con la comprobación realizada antes de aplicar cada idea.
